\section{Derivatives}

We start by defining differential operators in $p\geq1$ variables.
For a multi-index $i=(i_1,\dots,i_p)\in\N^p$, define the \emph{Hasse differential operator}
$\Delta_i$ by its effect on monomials:
\[
\Delta_i(x_1^{k_1}\cdots x_p^{k_p})=
\begin{cases}
\displaystyle\left(\prod_{j=1}^p {k_j\choose i_j}\right)
 x_1^{k_1-i_1}\cdots x_p^{k_p-i_p},& i_j\leq k_j\text{ for every }j,\\[3pt]
0,&\text{otherwise}.
\end{cases}
\]
We extend this operator linearly to all polynomials. It has order
$|i|=i_1+\cdots+i_p$. The definition works over $\Z$, $\Q$, or $\R$
(and over any commutative semiring by casting the binomial coefficients).
Over a $\Q$-algebra it agrees with
\[
\Delta_i=\frac{1}{i_1!\cdots i_p!}
 \left(\frac{\partial}{\partial x_1}\right)^{i_1}\cdots
 \left(\frac{\partial}{\partial x_p}\right)^{i_p}.
\]
When $i_j\leq k_j$, one may equivalently write
${k_j\choose i_j}={k_j\choose k_j-i_j}$.
Unless explicitly stated otherwise, a differential operator below means one
of these individual normalised operators, not an arbitrary linear combination.
This algebraic definition avoids the technicalities of Fr\'echet derivatives.

\lean{Represent polynomials by \texttt{MvPolynomial} and multi-indices
by finitely supported maps $\sigma\to_0\N$. Define each $\Delta_d$ as a
linear map and prove the coefficient identity
$[X^e]\Delta_dP=(\prod_j\binom{e_j+d_j}{d_j})[X^{e+d}]P$,
where $[X^e]$ denotes coefficient extraction; this shows that distinct
surviving monomials do not merge under differentiation. Use the univariate
